\documentclass[10pt,a4paper]{amsart}
\usepackage[margin=19mm]{geometry}
\usepackage{amsmath,amssymb,mathtools}
\usepackage[scr=rsfs,cal=euler]{mathalfa}
\usepackage{xfrac}
\usepackage[unicode,hidelinks,bookmarks=false]{hyperref}
\newcommand{\ie}{i.e.\ }
\newcommand{\cf}{cf.\ }
\newcommand{\resp}{respectively\ }
\newcommand{\Subsection}{\subsection}
\providecommand{\nobreakdash}{\nobreak\hbox{-}}
\setlength{\emergencystretch}{2em}
\multlinegap0pt
\usepackage{mathtools}
\usepackage[all]{xy}
\usepackage{amssymb}
\usepackage[shortlabels]{enumitem}
\newtheorem{prop}{Proposition}
\newtheorem{theorem}{Theorem}
\newtheorem{defin}{Definition}
\newtheorem{coro}{Corollary}
\newcommand{\Cat}{{\mathcal{C}\mspace{-2.mu}\it{at}}}
\newcommand{\Hf}[1]{\operatorname{H_{#1}}}
\newcommand{\Hotab}{\operatorname{\mathsf{Hotab}}}
\newcommand{\W}{\mathcal{W}}
\newcommand{\Wab}{\W^{\ab}}
\newcommand{\ab}{{\mathsf{ab}}}
\newcommand{\mdvirg}{\text{ ,}\hspace{10pt}}
\newcommand{\overbar}[1]{\mkern 1.5mu\overline{\mkern-1.5mu#1\mkern-1.5mu}\mkern 1.5mu}
\newcommand\pbox[1]{\zbox{\quad #1}} %
\newcommand{\pref}[1]{{\widehat{ #1 }}}
\newcommand{\prefab}[1]{\widehat{ #1 }^{\mathsf{ab}}}
\newcommand\zbox[1]{\makebox[0pt][l]{#1}}
\title{A Homotopical Dold-Kan correspondence for Joyal's category $\Theta$ and other test categories}
\author{L\'eo Hubert}
\date{Source: arXiv:2602.09674v1 (CC BY 4.0)}
\begin{document}
\maketitle

\noindent\emph{Source and licence.} A Homotopical Dold-Kan correspondence for Joyal's category $\Theta$ and other test categories. Version arXiv:2602.09674v1 (CC BY 4.0), distributed by arXiv under the Creative Commons Attribution 4.0 International licence, http://creativecommons.org/licenses/by/4.0/, as recorded in the arXiv licence metadata for this exact version and shown on the abstract page https://arxiv.org/abs/2602.09674v1. Full licence notice: \texttt{licenses/arxiv-2602.09674-CC-BY-4.0.txt}.\par\smallskip

\noindent\textit{Editorial scope.} Counted unit: the article's theorem thm:WhiteheadTestHomPseudoTest (source0.tex lines 2090--2093). The article's complete original proof of it is delivered with it (source0.tex lines 2094--2142, one \texttt{proof} environment of 49 lines). No mathematics is written, repaired or re-derived: the statement and the proof are the article's own lines, in the article's own notation, and no retained line is substituted or re-flowed.  Retained uncounted as the source-local context the proof needs: lines 623--630; lines 635--638; lines 688--692; lines 697--707; lines 721--731; lines 989--992; lines 1861--1864; lines 1964--1970; lines 2003--2017; lines 2058--2081.  Nothing was omitted from any retained span, so the package carries no omission record.  No retained line contains a citation, so the wrapper carries no bibliography and no bibliography entry.  No retained line contains a figure environment, an image-inclusion command or drawing code, so no image is delivered and the package declares no asset dependency.  Editorial formatting only, in addition: an AMS \texttt{amsart} preamble that restates the article's own declarations and macros used by the retained lines, namely the environments \texttt{prop}, \texttt{theorem}, \texttt{defin}, \texttt{coro} on separate counters, together with the standard packages those lines need.  The retained environments print in this document as Proposition 1, Theorem 1, Defin 1, Coro 1 in order of appearance, whereas the article prints the same items with the numbering of its own full document.  The complete native source of the article as distributed in the arXiv e-print for this exact version is bundled unchanged as \texttt{sources/194335/source0.tex}.
\begin{prop}[Grothendieck]\label{prop:totalementAspheriqueEquivalences}
Let $A$ be a small category. The following conditions are equivalent: 
\begin{enumerate}
\item $A$ is totally aspherical;
\item the diagonal functor $A \to A \times A$ is aspherical;
\item every aspherical presheaf on $A$ is locally aspherical.
\end{enumerate}
\end{prop}

\begin{prop}[Grothendieck]\label{totAspheriqueMorphismeAspherique}
If $A$ is a totally aspherical category and $u: A \to B$ is an aspherical
functor, then $B$ is also totally aspherical.
\end{prop}

\begin{prop}[Grothendieck]
\label{prop:testFaiblei_AAspherique}
Let $A$ be a small category. Then $A$ is a weak test category if and only if
$i_A$ is an aspherical functor.
\end{prop}

\begin{prop}[Grothendieck]\label{lemmeFoncteursAspheriquesTriCommutatif}
Let $u: A \to B$ be a functor between small categories, and $j: B \to \Cat$ a
functor such that for every object $b$ of $B$, the category $j(b)$ is
aspherical.
\begin{enumerate}
\item if $u$ is aspherical, then $ju: A \to \Cat$ is aspherical if
and only if $j$ is aspherical;
\item if $j$ is fully faithful and $ju: A \to \Cat$ is an aspherical functor, then
$u$ and $j$ are aspherical.
\end{enumerate}
\end{prop}

\begin{theorem}[Grothendieck]
\label{thm:EquivalencesFoncteursTestLocal}
Let $A$ be a small category and $i:A \to \Cat$ a functor such that for
every object $a$ of $A$, the category $i(a)$ is aspherical. The
following conditions are equivalent:
\begin{enumerate}
\item $i$ is a local test functor (and $A$ is a local test category);
\item  for every small aspherical category $C$, the presheaf $i^*(C)$
is locally aspherical.
\end{enumerate}
\end{theorem}

\begin{defin}\label{def:homPseudoTest}
A small category $A$ is a \emph{homologically pseudo-test category} if
the functor $\overbar{\Hf{A}}$ is an equivalence of categories.
\end{defin}

\begin{prop}\label{prop:aspheriqueSourceWhitehead}
Let $u: A \to B$ be an aspherical functor. If $A$ is a strong Whitehead
category, then $B$ is also a strong Whitehead category.
\end{prop}

\begin{theorem}\label{thm:modelStructureAbPresheaves}
Let $A$ be a strong Whitehead and local test category. The category
$\prefab{A}$ can be given the structure of a model category where weak
equivalences are the elements of~$\Wab_A$ and fibrations are the morphisms of
abelian presheaves whose underlying morphism of presheaves is a
fibration in the Grothendieck-Cisinski model structure on $\pref{A}$.
\end{theorem}

\begin{prop}\label{prop:foncteurAspheriqueEqQuillenAb}
Let $A$ be a strong Whitehead and local test category, $B$ a local test
category and $u: A \to B$ an aspherical functor, Then:
\begin{enumerate}
\item $B$ is a strong Whitehead category, and can be endowed with the model
structure defined in Theorem \ref{thm:modelStructureAbPresheaves};
\item the functors
\[
  (u^*)^\ab: \prefab{B} \to \prefab{A}\mdvirg u_*^\ab: \prefab{A} \to
  \prefab{B}
\]
forms a Quillen equivalence for the model structure defined in Theorem
\ref{thm:modelStructureAbPresheaves}.
\end{enumerate}
\end{prop}

\begin{coro}\label{coro:fonctAspheriqueTestLocalWhiteheadPsTestHomEq}
Let $A$ be a strong Whitehead and local test category, $B$ a local test
category, and~$u: A \to B$ an aspherical functor. Then $A$ is a
homologically pseudo-test category (\ref{def:homPseudoTest}) if and only if the same holds for~$B$.
\end{coro}
\begin{proof}
By Proposition \ref{prop:aspheriqueSourceWhitehead}, since $u$ is
aspherical and $A$ is a strong Whitehead category, then $B$ is a strong Whitehead
category. Thus, we can endow $\prefab{B}$ with the model category
structure described in Theorem \ref{thm:modelStructureAbPresheaves}.
Since aspherical morphisms are $\Wab_\infty$-aspherical, the diagram
\[
\xymatrix{
\Hotab_B \ar[rr]^{\overbar{(u^*)^\ab}} \ar[rd]_{\Hf{B}} && \Hotab_{A}
\ar[ld]^{\Hf{A}} \\
& \Hotab 
} 
\]
is commutative up to a natural isomorphism, and Proposition
\ref{prop:foncteurAspheriqueEqQuillenAb} 
shows that the horizontal arrow is an equivalence of categories. We
can then conclude that $\Hf{A}$ is an equivalence of categories if and
only if $\Hf{B}$ is an equivalence of categories.
\end{proof}

\begin{theorem}\label{thm:WhiteheadTestHomPseudoTest}
If $A$ is a strong Whitehead and test category, then $A$ is a
homologically pseudo-test category.
\end{theorem}

\begin{proof}
Let $A_0$ be a homologically pseudo-test, strong Whitehead and
strict test category (for example, we can take the category $\Delta$
of simplices). Let $B$ be the full subcategory of $\Cat$ whose objects
are the images by $i_A$ of objects of $A$ and images by $i_{A_0}$ of
objects of $A_0$. We then get the following commutative diagram
\[
\xymatrix{
& \Cat 
\\
A \ar[r]_{u} \ar[ru]^{i_A} &
B \ar@{^{(}->}[u]^{i} &
A_0 \ar[l]^{v} \ar[lu]_{i_{A_0}} 
} 
\]
in $\Cat$. We will show that $u$ and $v$ are aspherical functors, and
that $B$ is a strong Whitehead and local test category, before applying
Corollary \ref{coro:fonctAspheriqueTestLocalWhiteheadPsTestHomEq}.
Since $A_0$ is a local test category, the functor~$i_{A_0}$ is
aspherical by Proposition \ref{prop:testFaiblei_AAspherique}.
Moreover $i$ is full and faithful, so Proposition
\ref{lemmeFoncteursAspheriquesTriCommutatif} implies that $i$ and $v$ are
aspherical. Since $A$ is a local test category, the same argument shows that
$i_A$ and $u$ are also aspherical. Since $A_0$ is a strong Whitehead category
and $v$ is aspherical, Proposition \ref{prop:aspheriqueSourceWhitehead} implies
that $B$ is a strong Whitehead
category,

Let's show that $B$ is a local test category. 
Consider a small aspherical category $C$. Since $i$ is aspherical,
$i^*(C)$ is an aspherical presheaf on $B$ by Theorem
\ref{thm:EquivalencesFoncteursTestLocal}. But since $v$ is aspherical and $A_0$
is totally aspherical,
Proposition \ref{totAspheriqueMorphismeAspherique} implies that $B$ is also
totally aspherical. This implies by Proposition
\ref{prop:totalementAspheriqueEquivalences} that $i^*(C)$ is a locally
aspherical presheaf. By Theorem
\ref{thm:EquivalencesFoncteursTestLocal}, this implies that $B$ is a
local test category and that $i$ is a local test functor.
Finally, since $A$, $A_0$ and $B$ are strong Whitehead and local test
categories, and since there is a zig-zag of aspherical functors
\[
A \xrightarrow{u} B \xleftarrow{v} A_0 \pbox{,}
\]
we can conclude using Corollary
\ref{coro:fonctAspheriqueTestLocalWhiteheadPsTestHomEq} that $A$ is a
homologically pseudo-test category, since the same holds for
$A_0$.
\end{proof}

\end{document}
